在相同实际电价序列下，6月21日允许调整策略的费用为24865.86元，较不调整策略减少926.02元，且当日未发生紧急购电。12月21日允许调整后，紧急购电量由686.88 kWh降至554.17 kWh，但实际电费增加679.55元。经济性评价需综合原计划费、调整费与紧急购电费，不能仅依据紧急购电量。

两类策略的日初储电量分别由此前连续运行状态决定，指定展示日期不重新设定统一初值。因此，日表用于解释实际运行特征，策略总体比较仍采用全期费用，并单独校正期初与期末储电量差异。
